% partADC_6.tex — Phan 6: Adaptive Density Control minh hoa qua toy model 2D
% Chi chua cac khoi frame, khong documentclass/usepackage/begin{document}.

% --- 6_1_muc-tieu-va-gaussian-2d.png ---
\begin{frame}{Toy model 2D: mục tiêu và một Gaussian}
\begin{itemize}
    \footnotesize
    \item Ảnh mục tiêu $64\times 64$ px: một hình tròn phẳng (màu đều) + vùng chi tiết (dải mảnh 2 px, 5 chấm màu xen kẽ, ô vuông nhỏ).
    \item Mỗi Gaussian 2D: $G(x) = \exp\!\big(-\tfrac{1}{2}d^\top \Sigma^{-1} d\big)$, với $d = x - \mu$, $\Sigma$ tham số hoá bởi scale $s=(s_1,s_2)$ và góc xoay $\theta$.
    \item Footprint (vùng Gaussian được phép chạm) giới hạn ở $3\sigma$, tức $d^\top\Sigma^{-1}d < 9$ — ngoài phạm vi này chỉ còn màu nền.
    \item Render bằng alpha-compositing theo thứ tự chỉ số (độ sâu): $I(x) = \sum_i T_i\,\alpha_i G_i(x)\,c_i + T_{\text{end}}\cdot\text{bg}$, với $T_i = \prod_{j<i}(1-\alpha_j G_j)$.
\end{itemize}
\begin{center}
\includegraphics[width=0.75\linewidth]{6_1_muc-tieu-va-gaussian-2d.png}
\captionof{figure}{\footnotesize (a) Ảnh mục tiêu + lưới pixel; (b) một Gaussian $s=(7,3)$ px, $\theta=30°$ với các đường mức $1\sigma,2\sigma,3\sigma$; (c) kết quả render của Gaussian đó lên nền.}
\end{center}
\end{frame}

% --- 6_2_thieu-gaussian.png ---
\begin{frame}{Trường hợp thiếu Gaussian (under-reconstruction)}
\begin{itemize}
    \footnotesize
    \item Khởi tạo toy chỉ với 6 Gaussian nhỏ ($s=3$ px), rải rác không đủ để phủ vùng chi tiết (dải + chấm + ô vuông).
    \item Sai số $|I-T|$ (trung bình 3 kênh) lớn tại đúng những vùng không có Gaussian nào chạm tới.
    \item Gradient vị trí $-\nabla_\mu L$ tại mỗi Gaussian có độ lớn $\|\nabla_\mu\|$ vượt ngưỡng clone toy $\tau_{\text{grad}} = 4\times 10^{-5}$.
    \item Đây chính là tín hiệu để ADC quyết định \textbf{CLONE}: Gaussian di chuyển không đủ nhanh để tự phủ hết vùng thiếu, cần thêm bản sao.
\end{itemize}
\begin{center}
\includegraphics[width=0.75\linewidth]{6_2_thieu-gaussian.png}
\captionof{figure}{\footnotesize (a) Render với 6 Gaussian nhỏ; (b) mục tiêu; (c) bản đồ sai số; (d) mũi tên $-\nabla_\mu L$ và giá trị $\|\nabla_\mu\|\times 10^5$ tại từng Gaussian.}
\end{center}
\end{frame}

% --- 6_3_du-gaussian-qua-to.png ---
\begin{frame}{Trường hợp Gaussian dư / quá to (over-reconstruction)}
\begin{itemize}
    \footnotesize
    \item Một Gaussian duy nhất, kích thước lớn $s=(13,\ 6.5)$ px, cố phủ toàn bộ vùng chi tiết (dải + 5 chấm) cùng lúc.
    \item Sai số mờ đều trên toàn vùng, sai cả tại chấm lẫn khoảng trống — Gaussian quá to không tái tạo được chi tiết cục bộ.
    \item Ngưỡng phân loại ``to'': $\max_k s_k > S_{\text{big}} = 3.5$ px.
    \item Chỉ số quyết định split dùng \emph{tổng trị tuyệt đối} gradient theo pixel $\sum_x |g_x|$ (ngưỡng $\tau_{\text{abs}}=4\times 10^{-4}$), lớn hơn nhiều lần so với tổng có dấu $\|\sum_x g_x\|$ — vì các mũi tên gradient trong footprint kéo về nhiều phía trái ngược nhau.
\end{itemize}
\begin{center}
\includegraphics[width=0.75\linewidth]{6_3_du-gaussian-qua-to.png}
\captionof{figure}{\footnotesize (a) Render với 1 Gaussian to; (b) mục tiêu; (c) sai số mờ đều; (d) gradient theo từng pixel — tổng có dấu nhỏ nhưng tổng trị tuyệt đối lớn $\Rightarrow$ \textbf{SPLIT}.}
\end{center}
\end{frame}

% --- 6_4_clone-truoc-sau.png ---
\begin{frame}{CLONE: sao chép nguyên trạng Gaussian thiếu}
\begin{itemize}
    \footnotesize
    \item Trước clone: 1 Gaussian nhỏ có $\|\nabla_\mu\| \ge \tau$, cần phủ vùng rộng hơn kích thước hiện tại.
    \item CLONE tạo bản sao \emph{y hệt} (cùng $\mu,\Sigma,c,\alpha$); trạng thái Adam của bản mới được khởi tạo lại $m=v=0$ trong khi bản gốc giữ nguyên lịch sử — nên hai bản dịch chuyển khác tốc độ.
    \item Ngay sau clone hai bản trùng khít nhau; sau khoảng 120 bước Adam, quỹ đạo $\mu$ của bản gốc và bản sao tách rời để mỗi bản phủ một phần vùng cần thiết.
    \item Kết quả: loss L1 giảm rõ rệt sau khi hai bản tách ra và chia nhau vùng phủ.
\end{itemize}
\begin{center}
\includegraphics[width=0.75\linewidth]{6_4_clone-truoc-sau.png}
\captionof{figure}{\footnotesize (a) Trước clone; (b) ngay sau clone — 2 bản trùng nhau; (c) sau 120 bước — hai bản tách ra; (d) khoảng cách $\|\mu_{\text{gốc}}-\mu_{\text{sao}}\|$ theo bước, cùng L1 trước/sau.}
\end{center}
\end{frame}

% --- 6_5_split-truoc-sau.png ---
\begin{frame}{SPLIT: tách Gaussian quá to thành các con nhỏ hơn}
\begin{columns}[T]
\begin{column}{0.52\textwidth}
\scriptsize
\begin{itemize}
    \item $n_{\text{child}}{=}2$ con lấy mẫu từ hình dạng gốc: $\mu_{\text{con}}{=}\mu+R\epsilon^{(j)}$, $\epsilon^{(j)}{\sim}\mathcal{N}(0,\text{diag}(s^2))$.
    \item Scale mỗi con giảm còn $s/1.6$; Gaussian gốc bị \textbf{xoá}.
    \item Ngay sau split, loss \emph{xấu tạm thời}; sau 150 bước Adam mỗi con tự định vị lại.
    \item Kết quả cuối: L1 giảm, chi tiết cục bộ khôi phục tốt hơn 1 Gaussian to.
\end{itemize}
\end{column}
\begin{column}{0.48\textwidth}
\centering
\includegraphics[width=0.95\linewidth]{6_5_split-truoc-sau.png}
\captionof{figure}{\footnotesize Trên: vị trí/hình dạng trước–ngay sau–150 bước sau split. Dưới: render + L1.}
\end{column}
\end{columns}
\end{frame}

% --- 6_6_chuoi-thoi-gian.png ---
\begin{frame}{Chuỗi thời gian: toàn bộ vòng đời densify trên toy model}
\begin{itemize}
    \footnotesize
    \item Khởi tạo kiểu SfM-thu nhỏ: $N_0 = 12$ Gaussian, vị trí lấy mẫu ưu tiên vùng có gradient ảnh mục tiêu lớn (vùng chi tiết).
    \item Huấn luyện $T_{\text{total}}=1200$ vòng bằng Adam; ADC chạy mỗi $K=100$ vòng, trong cửa sổ densify $t\in[100, 900]$.
    \item Các mốc quan sát: $t = 0, 100, 300, 600, 1000$ — tại mỗi mốc ghi lại $N$ và L1 loss.
    \item Qua các mốc, Gaussian tự dịch chuyển, co giãn, nhân đôi (clone/split) và dần phủ kín ảnh mục tiêu, loss giảm liên tục.
\end{itemize}
\begin{center}
\includegraphics[width=0.85\linewidth]{6_6_chuoi-thoi-gian.png}
\captionof{figure}{\footnotesize Hàng trên: render tại từng mốc thời gian (N và L1 loss). Hàng dưới: ellipse $1.5\sigma$ của từng Gaussian chồng lên ảnh mục tiêu mờ.}
\end{center}
\end{frame}

% --- 6_7_chia-vung.png ---
\begin{frame}{Chia vùng: mỗi Gaussian ``phụ trách'' một lãnh thổ}
\begin{itemize}
    \footnotesize
    \item Với mỗi pixel, xác định Gaussian đóng góp lớn nhất theo trọng số $T_i \alpha_i G_i$ (đóng góp thực sự sau alpha-compositing).
    \item Pixel được tô theo màu của Gaussian ``thắng'' tại đó; màu trắng là nền (khi $T_{\text{end}}$ tại pixel đó lớn hơn mọi đóng góp).
    \item Đường viền đen là biên vật thể mục tiêu (đường mức gradient), dùng để đối chiếu vùng lãnh thổ với hình dạng thật.
    \item Theo thời gian, số lượng lãnh thổ tăng lên (do clone/split) và các lãnh thổ co lại, khớp sát hơn với biên vật thể và vùng chi tiết.
\end{itemize}
\begin{center}
\includegraphics[width=0.85\linewidth]{6_7_chia-vung.png}
\captionof{figure}{\footnotesize Bản đồ lãnh thổ tại các mốc thời gian: pixel tô theo Gaussian có $T_i\alpha_i G_i$ lớn nhất; đường đen = biên vật thể mục tiêu.}
\end{center}
\end{frame}

% --- 6_8_duong-cong-loss-N.png ---
\begin{frame}{Đường cong loss và số Gaussian $N(t)$}
\begin{columns}
\begin{column}{0.5\textwidth}
\begin{itemize}
    \footnotesize
    \item So sánh 3 kịch bản: (i) có ADC; (ii) không ADC, $N$ cố định bằng $N_0$; (iii) không ADC, $N$ cố định bằng $N$ cuối của ADC nhưng khởi tạo ngẫu nhiên ngay từ đầu.
    \item Vạch dọc đánh dấu các mốc ADC (mỗi $K_{\text{ADC}}=100$ vòng).
    \item $N(t)$ tăng theo bậc thang đúng tại các mốc ADC, chỉ trong cửa sổ densify $[100,900]$.
\end{itemize}
\end{column}
\begin{column}{0.5\textwidth}
\begin{itemize}
    \footnotesize
    \item Biểu đồ cột: số Gaussian bị clone / split / prune tại từng mốc ADC — cho thấy tỉ lệ thay đổi giữa các cơ chế qua thời gian.
    \item Có ADC luôn đạt L1 loss thấp hơn rõ rệt so với giữ nguyên $N_0$ ban đầu trong suốt quá trình.
\end{itemize}
\end{column}
\end{columns}
\begin{center}
\includegraphics[width=0.95\linewidth]{6_8_duong-cong-loss-N.png}
\captionof{figure}{\footnotesize (a) Loss (thang log, vạch dọc = mốc ADC); (b) $N(t)$ dạng bậc thang; (c) số clone/split/prune tại mỗi mốc ADC.}
\end{center}
\end{frame}

% --- 6_9_kich-thuoc-vs-chi-tiet.png ---
\begin{frame}{Kích thước Gaussian so với mức chi tiết tái tạo}
\begin{itemize}
    \footnotesize
    \item Phân bố kích thước $\max_k s_{i,k}$ (scale lớn nhất mỗi Gaussian) tại $t=0$ so với lúc kết thúc — Gaussian cuối tập trung nhiều hơn quanh và dưới ngưỡng toy $S_{\text{big}}=3.5$ px.
    \item Độ chi tiết vùng phủ của một Gaussian được đo bằng gradient ảnh mục tiêu trung bình trong footprint ($q<9$) của Gaussian đó.
    \item Quan hệ ngược chiều rõ rệt: Gaussian nằm ở vùng chi tiết cao có kích thước nhỏ hơn; Gaussian ở vùng phẳng (ví dụ hình tròn màu đều) có kích thước lớn hơn.
    \item Màu điểm biểu diễn opacity $\alpha$ — hầu hết Gaussian còn lại sau ADC có $\alpha$ cao, phản ánh chúng đóng góp thực chất vào ảnh cuối.
\end{itemize}
\begin{center}
\includegraphics[width=0.8\linewidth]{6_9_kich-thuoc-vs-chi-tiet.png}
\captionof{figure}{\footnotesize (a) Histogram kích thước Gaussian đầu vs cuối, vạch đỏ = ngưỡng $S_{\text{big}}$; (b) scatter độ chi tiết vs kích thước, màu = opacity $\alpha$.}
\end{center}
\end{frame}

% --- 6_10_khong-adc-vs-co-adc.png ---
\begin{frame}{So sánh trực tiếp: không ADC vs có ADC}
\begin{columns}[T]
\begin{column}{0.5\textwidth}
\scriptsize
\begin{itemize}
    \item Ba kịch bản: (1) không ADC, $N{=}N_0$; (2) không ADC, khởi tạo ngẫu nhiên $N{=}N_{\text{cuối}}$; (3) có ADC, $N_0{\to}N_{\text{cuối}}$ dần.
    \item (2) cùng số Gaussian (3) nhưng khởi tạo ngẫu nhiên không khớp tự nhiên vào vùng chi tiết.
    \item[$\Rightarrow$] ADC không chỉ tăng $N$ mà đặt đúng chỗ cần thiết.
\end{itemize}
\end{column}
\begin{column}{0.5\textwidth}
\centering
\includegraphics[width=0.85\linewidth]{6_10_khong-adc-vs-co-adc.png}
\captionof{figure}{\footnotesize 3 hàng (không ADC $N_0$; ngẫu nhiên; có ADC) $\times$ 3 cột.}
\end{column}
\end{columns}
\end{frame}

% --- 6_11_prune-opacity.png ---
\begin{frame}{PRUNE: loại bỏ Gaussian dư theo opacity}
\begin{itemize}
    \footnotesize
    \item Tại mốc $t=550$, opacity được reset: $\alpha \leftarrow \min(\alpha,\ \alpha_{\text{reset}}=0.02)$ — buộc mọi Gaussian phải ``chứng minh lại'' sự cần thiết của mình.
    \item Gaussian nào không phục hồi được opacity (không đóng góp giảm loss) sẽ bị đánh dấu để prune tại mốc ADC kế tiếp, theo ngưỡng $\alpha < \varepsilon_{\text{opa}} = 0.03$.
    \item So sánh render ngay trước và ngay sau prune: sai lệch tối đa $\max|\Delta I|$ rất nhỏ — ảnh gần như không đổi dù đã loại bỏ các Gaussian dư.
    \item Kết quả: mô hình gọn hơn (ít Gaussian hơn) mà chất lượng tái tạo giữ nguyên — đúng vai trò dọn dẹp của bước prune trong chu trình ADC.
\end{itemize}
\begin{center}
\includegraphics[width=0.8\linewidth]{6_11_prune-opacity.png}
\captionof{figure}{\footnotesize (a) Quỹ đạo opacity quanh mốc reset (đỏ = bị prune); (b) render ngay trước reset; (c) trước prune, ellipse đỏ = $\alpha<\varepsilon$; (d) sau prune — ảnh gần như không đổi.}
\end{center}
\end{frame}
